\section{Scope of the general conclusion}\label{sec:scope}
The main theorem is an equivalence for a compact continuous-response family with uniformly integrable physical returns. It is neither a declaration that every experiment has such returns nor a requirement that every slow boundary mode be uniformly controlled for every reward. Its exact obstruction is task compatibility at limiting ergodic projections. The algebraic theorem makes that obstruction decidable on finite algebraic instruments; the separate recurrence schedule gives an operational near-optimal family when the average objective is not continuous. These are distinct conclusions with different hypotheses.

The geometric conclusion concerns actual scored occupation. A finite acquired dimension gives a causal rate only after global candidate-valid covering, actual-law stability, a common executable initialization when needed, and state accounting have been established. Optimizing the controller requires the uniform lower witness and a fixed-size attaining exploration in Theorem~\ref{thm:geometry}. The four-call experiment verifies these optimized hypotheses directly for its geometric factor and derives its diagnostic lower against all history-dependent strategies. The queue realization does not claim the same dimension theorem for the entire workload posterior.

Physical regeneration is a substantive raw assumption. It does not mean observer overwrite, and it is not inferred from an arbitrary return decomposition. Common-program compactness here uses finite-product response domination or an explicit direct construction; general nondominated adaptive families remain outside that proof. The prediction quotient also has separate standard-Borel realization assumptions. Positivity of kernels alone supplies none of these extra properties.

Several natural problems remain. A necessary-and-sufficient recursive geometric classification beyond the cover-and-stability certificate is not proved. The efficient synthesis complexity of the global algebraic certification is not classified. Nor are the optimal joint costs of description length, workspace, physical computation time, numerical precision and simulator memory. Arbitrary unbounded stopping rules, nonregenerative infinite physical histories and unbounded quantum operators require further arguments. No total-variation estimate in this article is promoted to a large-deviation equivalence.

The present results isolate a common mechanism nevertheless: only the dynamically persistent directions visible to executable tasks need be paid for in a task-specific causal transfer. The marked approximate-corrector price quantifies that payment, acquired occupation identifies the spatial distortion, and an executable program must pay separately for preserving and updating both. This connects the original experimental foundation to a matching finite-resource law without replacing it by one realization.
